\section*{Bab \ref{ch:b1:magnetostatics} --- Magnetostatika: Medan, Gaya, dan Dipol}

\begin{solution}{exo:b1:magnetostatics:1}
$B = \mu_0I/2\pi r = 2 \times 10^{-7} \times 10/0.01 = \qty{2e-4}{T}$. Untuk \qty{1}{T}:
$I = 2\pi rB/\mu_0 = \qty{5e4}{A}$. Nilai Bumi tercapai pada $r = 2 \times 10^{-7} \times
100/5 \times 10^{-5} = \qty{0.4}{m}$.
\end{solution}

\begin{solution}{exo:b1:magnetostatics:2}
Pusatnya: $\mu_0NI/2R = 4\pi \times 10^{-7} \times 200/0.1 = \qty{2.5}{mT}$. Pada $z = R$:
$\times(R^2/2R^2)^{3/2} = 2^{-3/2}$: \qty{0.89}{mT}. Pada $z = 10R$: $\times(1/101)^{3/2}
= 10^{-3}$: \qty{2.5}{\micro T}, sudah menjadi hukum dipol $1/z^3$-nya.
\end{solution}

\begin{solution}{exo:b1:magnetostatics:3}
$B = \mu_0nI = 4\pi \times 10^{-7} \times 1000 \times 5 = \qty{6.3}{mT}$; di ujungnya,
separuhnya: \qty{3.1}{mT}. Kawat \qty{0.5}{mm}: $2000$ lilitan tiap meter, sehingga
\qty{12.6}{mT} pada \qty{5}{A}.
\end{solution}

\begin{solution}{exo:b1:magnetostatics:4}
$F = ILB = 10 \times 1 \times 0.1 = \qty{1.0}{N}$; pada \qty{30}{\degree}, dikalikan $\sin30^\circ$:
\qty{0.5}{N}. Motornya: $100 \times \qty{1}{N} \times \qty{0.05}{m} = \qty{5}{N.m}$ ---
sebuah motor industri yang kecil.
\end{solution}

\begin{solution}{exo:b1:magnetostatics:5}
Dengan $z = d\tan\alpha$: $B = \dfrac{\mu_0I}{4\pi}\displaystyle\int\dfrac{d\,\dd z}{(d^2 + z^2)^{3/2}}
= \dfrac{\mu_0I}{4\pi d}\displaystyle\int_{\alpha_1}^{\alpha_2}\cos\alpha\,\dd\alpha = \dfrac{\mu_0I}{4\pi d}
(\sin\alpha_2 - \sin\alpha_1)$; kawat tak hingga, $\alpha_{1,2} = \mp\pi/2$: $\mu_0I/2\pi d$.
Persegi: tiap sisinya pada $d = a/2$ yang terlihat di bawah $\pm45^\circ$ memberi $\mu_0I\sqrt2/2\pi a$;
empat sisinya: $B = 2\sqrt2\,\mu_0I/\pi a = 0.90\,\mu_0I/a$ (sebuah lingkaran ber-``jari-jari''
sama $a/2$ akan memberi $\mu_0I/a$).
\end{solution}

\begin{solution}{exo:b1:magnetostatics:6}
$B(z) = f(z - R/2) + f(z + R/2)$ dengan $f(u) = \mu_0NIR^2/2(R^2 + u^2)^{3/2}$.
Titik tengahnya: $2f(R/2) = \mu_0NIR^2/(5R^2/4)^{3/2} = (4/5)^{3/2}\mu_0NI/R$. $f'$
ganjil, sehingga $B'(0) = f'(-R/2) + f'(R/2) = 0$; $f''(u) \propto (4u^2 - R^2)(R^2 +
u^2)^{-7/2}$ lenyap pada $u = R/2$, sehingga $B''(0) = 2f''(R/2) = 0$: medannya
seragam sampai orde ketiga. Angkanya: $0.716 \times 4\pi \times 10^{-7} \times 100/0.15 =
\qty{0.60}{mT}$.
\end{solution}

\begin{solution}{exo:b1:magnetostatics:7}
Permukaannya: $\mu_0I/2\pi a = 2 \times 10^{-7} \times 100/0.002 = \qty{10}{mT}$; pada
\qty{1}{mm}, di dalam, $\times r/a$: \qty{5}{mT}; pada \qty{1}{cm}: \qty{2}{mT}. $B$ naik
linear sampai permukaannya, lalu meluruh sebanding $1/r$. Sesumbunya: di luar selongsongnya
arus terlingkupinya nol, $B = 0$.
\end{solution}

\begin{solution}{exo:b1:magnetostatics:8}
$2 \times 10^{-7}$~N/m. Rel dayanya: $2 \times 10^{-7} \times 10^8/0.1 = \qty{200}{N/m}$;
pada \qty{100}{kA}: $\qty{2e4}{N/m}$ --- dua ton tiap meter. Penopang
isolator pada panel hubung dirancang bagi arus gangguannya, bukan arus
kerjanya.
\end{solution}

\begin{solution}{exo:b1:magnetostatics:9}
$\Gamma = mB_h = \qty{2e-6}{N.m}$; $W = 2mB_h = \qty{4e-6}{J}$; $J\ddot\theta =
-mB_h\theta$: $T = 2\pi\sqrt{J/mB_h} = 2\pi\sqrt{10^{-7}/2 \times 10^{-6}} =
\qty{1.4}{s}$. Ukurlah waktu osilasinya: $B_h = 4\pi^2J/mT^2$ setelah $m$ dan $J$
diketahui (Gauss mengukur $m$ secara terpisah dari simpangan yang diberikannya
pada jarum kedua).
\end{solution}

\begin{solution}{exo:b1:magnetostatics:10}
(a) $\mu_0IR^2/2z^3 = \mu_0(I\pi R^2)/2\pi z^3 = \mu_0m/2\pi z^3 = (\mu_0/4\pi)\,2m/z^3$:
yaitu $2p/4\pi\varepsilon_0z^3$ listriknya dengan $1/\varepsilon_0 \to \mu_0$. (b) $m = 2\pi R_E^3B/
\mu_0 = 2\pi \times 2.58 \times 10^{20} \times 6 \times 10^{-5}/1.26 \times 10^{-6} = \qty{7.7e22}{A.m^2}$.
(c) $I = m/\pi R^2 = 7.7 \times 10^{22}/(\pi \times 9 \times 10^{12}) = \qty{2.7e9}{A}$. (d)
Khatulistiwa ($\theta = 90^\circ$): $B_\theta = \mu_0m/4\pi R_E^3 = \qty{3e-5}{T}$, separuh
nilai kutubnya. Pada $\theta = 45^\circ$: $\tan(\text{celup}) = B_r/B_\theta = 2\cot\theta = 2$,
sudut celupnya $= \qty{63}{\degree}$ --- medannya menukik curam ke dalam tanah pada
lintang menengah.
\end{solution}

\begin{solution}{exo:b1:magnetostatics:11}
$\mu_0nI = 4\pi \times 10^{-7} \times 2000 \times 3 = \qty{7.5}{mT}$. Pusatnya: $\tan\alpha_1 =
R/(L/2) = 0.2$, $\cos\alpha_1 = -\cos\alpha_2 = 0.981$: $B = 0.981\,\mu_0nI =
\qty{7.4}{mT}$ ($98\%$). Ujungnya: $\cos\alpha_1 = 0$, $\tan\alpha_2 = R/L$, $\cos\alpha_2 =
-0.995$: $B = 0.497\,\mu_0nI = \qty{3.7}{mT}$ ($50\%$). Di luar pada $x$ melampaui
ujungnya, kedua ujungnya sesisi: $B = \tfrac12\mu_0nI(\cos\alpha_1 - \cos\alpha_2)
\approx \tfrac12\mu_0nI\,\tfrac{R^2}{2}[1/x^2 - 1/(x + L)^2]$; $1\%$ nilai pusatnya
tercapai pada $x \approx \qty{10}{cm}$, yaitu setengah panjang di luar. $B(z)$: puncak datar sepanjang
sebagian besar panjangnya, turun menjadi separuh di ujungnya dan menjadi beberapa persen
sejauh satu panjang.
\end{solution}

\begin{solution}{exo:b1:magnetostatics:12}
(a) Medan gulungannya sendiri melompat dari $0$ ke $B$ melintasinya;
medan yang bekerja padanya adalah rata-ratanya, $B/2$ (lembarnya tak dapat mendorong
dirinya sendiri). Gaya pada sebuah unsur gulungan sepanjang $\dd l$ dan berlebar $\dd z$
yang mengalirkan $K\,\dd z$: $K\,\dd z\,\dd l\,B/2$, radial ke luar (periksalah dengan $I\,\dd\vect l
\wedge\vect B$): tekanan $p = KB/2 = B^2/2\mu_0$ karena $K = nI = B/\mu_0$.
(b) \qty{1}{T}: \qty{4e5}{Pa} = \qty{4}{bar}; \qty{10}{T}: \qty{400}{bar}; \qty{45}{T}:
\qty{8e8}{Pa} = \qty{8000}{bar}, dekat dengan \qty{1}{GPa} milik baja. (c)
Gulungannya harus menahan tekanan yang tumbuh sebanding $B^2$: di atas kira-kira \qty{40}{T}
tak ada bahan dan tak ada skema pendinginan yang sanggup --- batas magnet tunak
bersifat mekanis (dan, bagi superkonduktor, medan kritisnya).
\end{solution}

\begin{solution}{pb:b1:magnetostatics:1}
\textbf{1.} Setiap bidang yang memuat sumbunya adalah \omterm{prop:b1:electrostatics-gauss:symmetry}{bidang simetri}
arusnya: $\vect B$ tegak lurus terhadapnya, ortoradial; invariansi
menyusuri dan mengelilingi sumbunya: $\vect B = B(r)\,\vect e_\theta$.

\textbf{2.} Lingkaran berjari-jari $r$: $2\pi rB = \mu_0I_{\text{dalam}}$. Terasnya:
$I_{\text{dalam}} = Ir^2/a^2$, $B = \mu_0Ir/2\pi a^2$; di antaranya: $\mu_0I/2\pi r$;
di luar ($r > c$): $I - I = 0$, $B = 0$.

\textbf{3.} $B(a) = 2 \times 10^{-7}/5 \times 10^{-4} = \qty{4.0e-4}{T}$; $B(b) = 2 \times
10^{-7}/2.5 \times 10^{-3} = \qty{8.0e-5}{T}$.

\textbf{4.} Di dalam tabungnya tak ada arus yang terlingkupi: $B = 0$ di sana alih-alih
kenaikan linearnya; di tempat lain tak ada yang berubah --- medan di luar
sebuah arus silinder hanya bergantung pada arus totalnya.

\textbf{5.} $I_{\text{dalam}} = I - I(r^2 - b^2)/(c^2 - b^2) = I(c^2 - r^2)/(c^2 -
b^2)$: $B = \dfrac{\mu_0I}{2\pi r}\,\dfrac{c^2 - r^2}{c^2 - b^2}$, sama dengan $\mu_0I/2\pi b$ di
$b$ dan $0$ di $c$: kontinu. Sketsanya: linear sampai \qty{4e-4}{T} pada
\qty{0.5}{mm}, lalu $1/r$ turun ke \qty{8e-5}{T} pada \qty{2.5}{mm}, turun menjadi nol pada
\qty{3}{mm}, nol di luarnya.

\textbf{6.} Arus selongsongnya berlawanan arah dengan arus terasnya: tolakan,
yaitu tekanan ke luar. Arus permukaan $K = I/2\pi b = \qty{64}{A/m}$; medan rata-ratanya
$\tfrac12(8 \times 10^{-5} + 0) = \qty{4e-5}{T}$: $p = KB = \qty{2.5e-3}{Pa}$. Pada
\qty{10}{kA}: $\times10^8$, \qty{2.5e5}{Pa} = \qty{2.5}{bar} --- kabel daya
denyut dibuat untuk menahannya.

\textbf{7.} Tak ada medan di luar: kabelnya tak memancar dan tak memungut
gangguan magnetik dari tetangganya. Sepasang kawat: kedua
medannya hampir saling meniadakan, menyisakan $\approx \mu_0Id/2\pi D^2 = 2 \times 10^{-7} \times 2 \times
10^{-3}/10^{-2} = \qty{4e-8}{T}$ pada \qty{10}{cm} --- kecil, tetapi sesumbunya memberi
nol.

\textbf{8.} $D = \sqrt{h^2 + d^2/4} = \qty{10.0}{m}$: masing-masing memberi $\mu_0I/2\pi D =
\qty{1.0e-5}{T}$, tegak lurus pada garis dari penghantarnya. Karena
arusnya berlawanan, komponen yang sejajar dengan garis penghubung
kedua penghantarnya saling meniadakan dan komponen tegaknya berjumlah: $B = 2 \times 10^{-5} \times
(d/2)/D = \qty{1.0e-6}{T}$, tegak. Lima puluh kali di bawah medan Bumi,
seratus kali di bawah pedomannya.

\textbf{9.} Untuk $D \gg d$, $B \approx \dfrac{\mu_0I}{2\pi}\Bigl(\dfrac1{D_1} - \dfrac1{D_2}\Bigr)
\sim \dfrac{\mu_0I}{2\pi}\dfrac{d}{D^2}$: kedua medan $1/D$-nya saling meniadakan sampai orde
pertama dan hanya selisihnya, yang berukuran nisbi $d/D$, yang bertahan --- sebuah
``dipol dua kawat''.

\textbf{10.} Penghantarnya sendiri pada \qty{1}{m}: \qty{1.0e-4}{T} mendatar; yang lain
pada \qty{1.41}{m}: \qty{7.1e-5}{T} pada \qty{45}{\degree}, berlawanan: resultannya
$(5 \times 10^{-5}, 5 \times 10^{-5})$, \qty{7e-5}{T} --- \qty{70}{\micro T}, di bawah
pedoman umum dan jauh di bawah pedoman kerjanya.

\textbf{11.} $F/\ell = \mu_0I^2/2\pi d = 2 \times 10^{-7} \times 2.5 \times 10^5 =
\qty{0.05}{N/m}$, bersifat menolak; pada \qty{20}{kA}: \qty{80}{N/m}.

\textbf{12.} \qty{15}{N} dalam keadaan biasa; \qty{24}{kN} dalam korsleting, delapan kali
berat rentangnya \qty{2.9}{kN}: penghantarnya terlempar berpisah, berayun
dan mungkin berbenturan ketika jatuh kembali --- karena itu ada pemisah dan pemutus yang
membuka dalam beberapa daur.

\textbf{13.} $I\cos\omega t + I\cos(\omega t - 2\pi/3) + I\cos(\omega t + 2\pi/3) =
I\cos\omega t\,(1 + 2\cos2\pi/3) = 0$: ketiga \omterm{def:b1:sinusoidal-impedance:complex}{fasornya} berjumlah nol. Medan
$1/D$-nya saling meniadakan seperti pada sesumbunya; medan jauhnya meluruh sebanding $1/D^2$ atau
lebih cepat.

\textbf{14.} Medan salurannya di sana tegak; sebuah kompas menanggapi
komponen mendatarnya saja: tak ada simpangan (sudut celupnya berubah sebesar
$10^{-6}/(2 \times 10^{-5})$, yaitu tiga derajat, yang tak dapat ditunjukkan kompasnya).

\textbf{15.} Sisi $a$ (sejajar dengan sumbunya): $\vect B$ yang radial
tegak lurus terhadapnya, gayanya $NIaB$ pada masing-masing, tegak lurus pada
bidang kumparannya, berlawanan pada kedua sisinya: sebuah kopel. Sisi $b$: $\dd\vect l$
radial, sejajar dengan $\vect B$: tak ada gaya.

\textbf{16.} $\Gamma = 2 \times NIaB \times b/2 = NIabB = NISB = 0.024\,I$ (N\,m, $I$
dalam A). Medan radialnya selalu di dalam bidang kumparannya dan tegak lurus
pada sisi $a$-nya: gayanya selalu normal terhadap bidangnya dan \omterm{def:b1:angular-momentum:moment}{lengan
gayanya} $b/2$ tak pernah berubah.

\textbf{17.} $C\theta = NISB$: $\theta = 2400\,I$ rad, yaitu \qty{2.4}{rad/mA}; skala
penuh \qty{1.2}{rad} pada $I = \qty{0.50}{mA}$.

\textbf{18.} Panjangnya $200 \times 2(0.02 + 0.03) = \qty{20}{m}$, penampangnya $\pi(5 \times
10^{-5})^2 = \qty{7.9e-9}{m^2}$: $R = 1.7 \times 10^{-8} \times 20/7.9 \times 10^{-9} =
\qty{43}{\ohm}$; $U = 43 \times 5 \times 10^{-4} = \qty{22}{mV}$; $P = UI = \qty{11}{\micro W}$.

\textbf{19.} $J\ddot\theta = -C\theta + NISB\,I$: $\omega_0 = \sqrt{C/J} = \sqrt{10^{-5}/6.7
\times 10^{-8}} = \qty{12}{rad/s}$, $T = \qty{0.51}{s}$.

\textbf{20.} $J\ddot\theta + \alpha\dot\theta + C\theta = NISB\,I$: kritis untuk
$\alpha = 2\sqrt{JC} = 2\sqrt{6.7 \times 10^{-13}} = \qty{1.6e-6}{N.m.s}$; jarumnya
lalu mencapai bacaannya paling cepat tanpa melampaui.

\textbf{21.} Voltmeter: $R_s = 10/5 \times 10^{-4} - 43 \approx \qty{20}{k\ohm}$.
Amperemeter: pirauannya mengalirkan \qty{0.9995}{A} di bawah \qty{22}{mV}: $R_{\text{pir}} =
\qty{22}{m\ohm}$.

\textbf{22.} $m = NIS = 200 \times 5 \times 10^{-4} \times 6 \times 10^{-4} = \qty{6e-5}{A.m^2}$;
$-mB = \qty{-1.2e-5}{J}$; pegasnya $\tfrac12C\theta^2 = \tfrac12 \times 10^{-5} \times 1.44 =
\qty{7.2e-6}{J}$: ordenya sama --- pegasnya menyimpan apa yang dikerjakan \omterm{def:b1:angular-momentum:moment}{momen
gaya} magnetiknya selama kumparannya berputar.

\textbf{23.} Dalam medan seragam \omterm{def:b1:angular-momentum:moment}{momen gayanya} $mB\cos\theta$ (maksimum ketika
bidangnya memuat $\vect B$, nol ketika normalnya menyusurinya): $C\theta =
NISB\,I\cos\theta$ --- linear hanya bagi $\theta$ kecil, mampat di
ujung atas skalanya.

\textbf{24.} \qty{50}{Hz} terhadap $\omega_0/2\pi = \qty{2}{Hz}$: kumparannya tak dapat
mengikuti; ia duduk pada rata-rata arusnya, yaitu nol, dengan getar
beramplitudo $\approx \theta_{\text{statik}}(\omega_0/\omega)^2 = 1.2 \times (12/314)^2 \approx
\qty{2e-3}{rad}$, tak terlihat. Alat ukur kumparan putar membaca rata-ratanya --- bagi
arus bolak-balik ia memerlukan penyearah.

\textbf{25.} \qty{4e-4}{T} pada teras sebuah kabel sesumbu \qty{1}{A} dan nol di luarnya
(\omterm{thm:b1:magnetostatics:ampere}{hukum Ampère}); \qty{1}{\micro T} di bawah saluran \qty{500}{A} dan \qty{0.05}{N/m}
antara penghantarnya (Biot--Savart ditambah Laplace); \qty{2.4}{rad/mA} bagi
alat ukurnya (kopel Laplace melawan sebuah pegas).
\end{solution}
